\documentclass[11pt]{article}
\usepackage[]{acl}
\usepackage{times}
\usepackage{latexsym}
\usepackage[T1]{fontenc}
\usepackage[utf8]{inputenc}
\usepackage{microtype}
\usepackage{graphicx}
\usepackage{booktabs}
\usepackage{amsmath}
\usepackage{multirow}

\title{Operating Point, Not Capacity: Calibration and Orthogonal-Feature
Routing for Robust and Interpretable AI-Generated Text Detection}

\author{Hanzhong Zhang \\
  Fudan University \\
  \texttt{24307140093} \\}

\begin{document}
\maketitle

\begin{abstract}
We revisit machine-generated text detection on SemEval-2024 Task~8 Subtask~A
(English) under its realistic distribution shift: the test set contains
generators and a domain never seen at training time. A fine-tuned
RoBERTa-base reaches test AUROC $0.944$ yet only $0.619$ macro-F1 at the
default $0.5$ threshold---a $0.27$ gap we trace not to representation
capacity but to a \emph{mis-placed operating point} under shift. Temperature
scaling plus a threshold chosen on an in-distribution calibration set
recovers macro-F1 to $0.886$ (oracle $0.888$) without retraining. We further
show that a single global threshold is structurally unable to catch the
human-like \texttt{bloomz} generator (recall $0.096$), because its decision
margin overlaps the human upper tail. We propose a \emph{region-aware
two-stage detector}: confident-margin documents are decided directly, while
an ambiguous band is resolved by a logistic classifier over orthogonal
hand-crafted features (lexical diversity, repetition). Under the strict
fit-on-dev / predict-test-once protocol this reaches macro-F1 $0.857$,
surpassing the best single-model baseline ($0.835$); the in-distribution
ceiling is $0.945$. The orthogonal stage lifts \texttt{bloomz} recall from
$0.10$ to $0.70$. We pair the detector with a plagiarism-checker-style
interpretable report (occlusion-based sentence attribution, a
human$\leftrightarrow$AI discriminant axis for each linguistic feature) and a
closed-loop de-AI robustness evaluation. A negative result---a hierarchical
sentence-aware extension that overfits the single-generator dev set
(test AUROC $0.944\!\rightarrow\!0.672$)---reinforces our thesis: the
bottleneck is calibration and shift, not model capacity.
\end{abstract}

\section{Introduction}

Detecting machine-generated text is increasingly framed as a problem of
\emph{generalization}: a deployed detector must flag text from generators and
domains it never saw during training. SemEval-2024 Task~8 Subtask~A
\citep{wang2024semeval} operationalizes exactly this---its English test split
introduces two unseen generators (GPT-4, \texttt{bloomz}) and a single unseen
domain (\texttt{outfox}), while the training data covers four other generators
across five domains.

A standard fine-tuned RoBERTa-base \citep{liu2019roberta} on this setup
exhibits a striking contradiction: test AUROC is $0.944$, but macro-F1 at the
conventional $0.5$ threshold is only $0.619$. High AUROC with low accuracy is
almost impossible unless the \emph{operating point} is wrong. We confirm this:
the model still ranks AI above human well under shift, but its absolute
probabilities are systematically displaced, so a fixed threshold collapses
(it falsely flags $11{,}147$ of $16{,}272$ human texts).

This observation reframes the improvement problem. The community reflex is to
add capacity---hierarchical encoders, multi-task heads, auxiliary
supervision. We show this reflex backfires here (Section~\ref{sec:hsad}). The
actual levers are (i) \textbf{calibration and threshold adaptation} and (ii) a
targeted fix for the one generator a global threshold cannot catch. Our
contributions:

\begin{itemize}\itemsep2pt
\item A diagnosis isolating the $0.27$ macro-F1 gap to operating point, not
representation: temperature scaling plus an in-distribution threshold recovers
$0.619\!\rightarrow\!0.886$ (oracle $0.888$) with the model frozen
(Section~\ref{sec:diag}).
\item \textbf{Region-aware two-stage detection}: a confident-margin fast path
plus an ambiguous-band classifier over \emph{orthogonal} hand-crafted features.
Under the strict fit-on-dev / test-once protocol it reaches macro-F1 $0.857$,
beating the strongest single-model baseline ($0.835$), and lifts the
human-like \texttt{bloomz} generator's recall from $0.10$ to $0.70$
(Section~\ref{sec:method}).
\item An interpretable, plagiarism-checker-style report: occlusion-based
sentence attribution that is coherent with the document score, and a
per-feature human$\leftrightarrow$AI discriminant axis (Section~\ref{sec:viz}).
\item A closed-loop de-AI robustness evaluation under cross-model recursive
paraphrase attacks, revealing that the orthogonal-feature stage is
simultaneously the source of \texttt{bloomz} recovery and the principal attack
surface (Section~\ref{sec:robust}).
\item A documented negative result (HSAD) showing that added hierarchical
capacity overfits the single-generator dev set and destroys ranking ability
(Section~\ref{sec:hsad}).
\end{itemize}

\section{Related Work}

\paragraph{Machine-generated text detection.} Approaches divide into
zero-shot statistical detectors, e.g.\ DetectGPT's probability-curvature test
\citep{mitchell2023detectgpt}; supervised fine-tuned classifiers, which
dominate leaderboards but generalize poorly off-distribution; and
watermarking \citep{kirchenbauer2023watermark}, which requires generator
cooperation. Surveys \citep{wu2025survey} note that perplexity and burstiness
remain the interpretable backbone signals, which we surface in our report.

\paragraph{Robustness to paraphrase.} \citet{krishna2023paraphrasing} show
paraphrasing evades most detectors and propose retrieval defenses;
recursive/iterative paraphrasing is the strongest attack. We reproduce this
under a cross-model setting where strong instruction models rewrite text that
our detector then re-scores.

\paragraph{Calibration.} Temperature scaling \citep{guo2017calibration} is a
standard post-hoc fix; we show that under distribution shift its benefit is
conditional on having an in-distribution calibration set, which the official
dev split is not.

\section{Setup and Diagnosis}
\label{sec:diag}

\paragraph{Data.} SemEval-2024 Task~8 Subtask~A (English): train $119{,}756$
docs (human + \texttt{davinci}, \texttt{chatGPT}, \texttt{dolly},
\texttt{cohere}; domains reddit/wikihow/arxiv/wikipedia/peerread), dev
$5{,}000$ (only \texttt{bloomz} as the AI generator, $5$ domains), test
$34{,}272$ (six generators incl.\ unseen GPT-4 and \texttt{bloomz}; single
\texttt{outfox} domain). The dev and test splits thus differ on \emph{both}
axes---generator and domain.

\paragraph{The AUROC/accuracy contradiction.} Table~\ref{tab:ceiling} shows
fine-tuned RoBERTa-base reaches test AUROC $0.944$ but macro-F1 $0.619$ at
threshold $0.5$. Per-group analysis reveals near-perfect recall on AI groups
but $31\%$ accuracy on human---the threshold is displaced, not the ranking.

\paragraph{Dev cannot calibrate test.} The dev-optimal margin threshold
($-5.86$) and the test oracle threshold ($+12.0$) point in \emph{opposite}
directions; applying the dev threshold to test makes accuracy \emph{worse}
($0.642 < 0.671$). Because dev contains only \texttt{bloomz}, a threshold
tuned on it does not transfer to GPT-4/chatGPT-style text.

\paragraph{Operating-point ceiling.} To estimate how much a frozen model can
be fixed \emph{given} an in-distribution calibration set, we stratify test
$50/50$ by $\text{generator}\times\text{label}$, fit temperature and a
macro-F1-optimal margin threshold on one half, and score the disjoint half
($5$ seeds). Table~\ref{tab:ceiling} reports the result: macro-F1
$0.619\!\rightarrow\!0.886$ (oracle $0.888$), confirming the gap is operating
point, not capacity. TF-IDF + logistic regression is already well-calibrated
(ECE $0.03$) and barely moves, explaining why this lexical baseline
\emph{outperforms} raw RoBERTa at the default threshold.

\begin{table}[t]\centering\small
\begin{tabular}{lccc}
\toprule
Model & base F1 & calib.\ F1 & AUROC \\
\midrule
RoBERTa single   & 0.619 & \textbf{0.886} & 0.944 \\
RoBERTa chunked  & 0.677 & 0.880 & 0.958 \\
TF-IDF + LR      & 0.828 & 0.830 & 0.897 \\
Fusion LGBM      & 0.835 & 0.882 & 0.930 \\
\bottomrule
\end{tabular}
\caption{Test macro-F1 at threshold $0.5$ (base) vs.\ after temperature +
in-distribution threshold (calib., mean over 5 seeds). The frozen RoBERTa
recovers $+0.27$; AUROC is threshold-invariant.}
\label{tab:ceiling}
\end{table}

\paragraph{Why \texttt{bloomz} is hard.} After calibration, per-generator AI
recall is $\geq0.99$ for seen-family generators and $0.89$ for unseen GPT-4,
but only $0.096$ for \texttt{bloomz}. In margin space, \texttt{bloomz}
(median $10.9$, p90 $12.0$) overlaps the human upper tail (p75 $11.2$, p90
$11.9$), while clean AI sits $\geq12.1$. A single global threshold therefore
trades human specificity against \texttt{bloomz} recall: at margin $12.0$,
specificity $0.95$ but \texttt{bloomz} recall $0.10$; lowering to $11.0$
raises \texttt{bloomz} to $0.48$ but drops specificity to $0.72$. The two are
in direct conflict for any single threshold.

\section{Region-Aware Two-Stage Detection}
\label{sec:method}

The conflict above is not a ranking failure but a \emph{single-threshold}
failure: \texttt{bloomz} is separable from human in \emph{orthogonal}
hand-crafted features (type-token ratio $0.83$ vs.\ $0.43$; repeated-bigram
ratio $0.01$ vs.\ $0.15$) even where the RoBERTa margin cannot separate them.
We exploit this with a two-stage rule. Let $m=\text{logit}_{\text{AI}}-
\text{logit}_{\text{human}}$ be the (temperature-scaled) document margin:
\begin{equation}
\hat{y}=
\begin{cases}
\text{AI} & m \geq t_{\text{high}}\\
\text{human} & m < t_{\text{low}}\\
g_\phi(\mathbf{f}) & t_{\text{low}}\leq m < t_{\text{high}}
\end{cases}
\end{equation}
where $g_\phi$ is a logistic classifier over six orthogonal features
(TTR, repeated bi/tri-gram ratios, avg.\ word length, length, punctuation
density). All of $t_{\text{low}}, t_{\text{high}}, \phi$ and the temperature
are fit on the dev split only (the deployable P1 protocol). The learned
$g_\phi$ is dominated by TTR (coefficient $16.3$), matching the diagnosis.

\paragraph{Protocols.} We report three numbers to separate effects honestly:
\textbf{P1} (deployable) fits everything on dev and predicts test once;
\textbf{P2} (in-distribution ceiling) uses a nested split of test for honest
band selection; \textbf{P0} is the optimistic in-test variant.
Table~\ref{tab:method} shows P1 macro-F1 $0.857$---\emph{above} the strongest
single-model baseline (Fusion LGBM, $0.835$, itself test-tuned)---and a P2
ceiling of $0.945$. Crucially the region-aware rule lifts \texttt{bloomz}
recall to $0.70$ (P2) / $1.00$ (P1) while keeping human specificity high, and
the P0$\approx$P2 gap of $0.0016$ shows the band selection does not overfit.

\begin{table}[t]\centering\small
\begin{tabular}{llcc}
\toprule
Protocol & rule & macro-F1 & note \\
\midrule
P1 deploy & single-thr & 0.572 & dev$\to$test \\
P1 deploy & region-aware & \textbf{0.857} & dev$\to$test \\
P2 ceiling & region-aware & 0.945 & in-dist \\
\midrule
\multicolumn{2}{l}{Fusion LGBM (prior best)} & 0.835 & test-tuned \\
\bottomrule
\end{tabular}
\caption{Region-aware vs.\ single-threshold. P1 is the only leakage-free test
number; it beats the prior best single model. Single-threshold under the same
P1 protocol collapses to $0.572$.}
\label{tab:method}
\end{table}

\section{Interpretable Report}
\label{sec:viz}

We render each document as a plagiarism-checker-style report with three
coherent layers. (1) The headline is the \emph{real} temperature-scaled
$P(\text{AI})$, not a discretized label, so it varies continuously and never
contradicts the model's confidence; the region-aware verdict and its decision
region are shown separately. (2) Sentence highlighting uses
\emph{occlusion}: a sentence's contribution is
$P(\text{AI}\mid\text{full})-P(\text{AI}\mid\text{without sentence})$, so the
highlight is by construction consistent with the document score (avoiding the
common artifact where independently-scored sentences all look benign yet the
document is flagged). (3) Each linguistic feature is placed on a
human$\leftrightarrow$AI \emph{discriminant axis} via a Gaussian
log-likelihood ratio,
$\text{disc}=\tanh\!\big(\tfrac{1}{2}[\log\mathcal{N}(x;\mu_{\text{AI}},
\sigma_{\text{AI}})-\log\mathcal{N}(x;\mu_{\text{H}},\sigma_{\text{H}})]\big)
\in[-1,1]$, so position encodes how decisively the feature leans AI, using
both class means and spreads. We also surface perplexity (GPT-2) and
burstiness as the classic interpretable signals.

\section{De-AI Robustness Evaluation}
\label{sec:robust}

We frame ``de-AI rewriting'' as a detector-robustness probe: an attacker
model rewrites a high-AI document using the detector's own feature diagnosis,
and we re-score it, over $K$ recursive rounds. Using an OpenAI-compatible
gateway we compare three attackers (\texttt{claude-sonnet-4-6},
\texttt{gpt-5.2}, \texttt{qwen3-max}), $12$ high-AI documents each. A single
rewrite barely moves the detector
($P(\text{AI})\,0.906\!\rightarrow\!0.90$); but recursive rewriting with a
strong attacker breaks it. Table~\ref{tab:robust} shows Claude flips $83\%$ of
documents to human by round~2 ($P(\text{AI})$ dropping
$0.906\!\rightarrow\!0.733$, in one case driving the margin from $+12.2$ to
$-7.3$), whereas weaker attackers flip only $25$--$33\%$ and barely move
$P(\text{AI})$. Attacker strength is the dominant factor.

\begin{table}[t]\centering\small
\begin{tabular}{lccc}
\toprule
Attacker & flip @r1 & flip @r2 & $P(\text{AI})$ r0$\to$r2 \\
\midrule
claude-sonnet-4-6 & 67\% & \textbf{83\%} & 0.906$\to$0.733 \\
qwen3-max         & 17\% & 33\% & 0.906$\to$0.884 \\
gpt-5.2           & 8\%  & 25\% & 0.906$\to$0.884 \\
\bottomrule
\end{tabular}
\caption{Recursive cross-model paraphrase attack ($12$ docs/attacker, $K{=}2$).
Flip = fraction of documents the detector re-labels human. A strong attacker
(Claude) breaks the detector; weak attackers largely do not.}
\label{tab:robust}
\end{table}

\paragraph{Where the attack lands.} The flip is driven by our
\emph{orthogonal-feature stage}, not RoBERTa: rewriting scrambles surface
features (TTR, contractions), pushing the margin into the ambiguous band where
$g_\phi$ then flips to human, while RoBERTa's deep $P(\text{AI})$ moves little
($0.906\!\rightarrow\!0.83$). Thus the orthogonal stage is \emph{both} the
mechanism that recovers \texttt{bloomz} and the principal attack surface---a
direct trade-off between human-like-text recall and paraphrase robustness.

\section{Negative Result: Hierarchical Overfitting}
\label{sec:hsad}

We also tested the capacity reflex: HSAD, which pools RoBERTa token states
into sentence vectors, adds a 2-layer cross-sentence Transformer and dual
sentence/document heads (bottom 8 RoBERTa layers frozen, multi-task with
sentence pseudo-labels). HSAD reaches dev accuracy $0.843$ (vs.\ $0.834$ for
plain RoBERTa) but its test AUROC \emph{collapses} from $0.944$ to $0.672$.
The added degrees of freedom overfit the single-generator (\texttt{bloomz})
dev distribution and fail to generalize to unseen generators, destroying the
ranking ability the backbone had. This negative result supports our central
claim: the bottleneck is operating point and shift, not capacity, so the
right fix is lightweight calibration plus orthogonal routing, not a heavier
model.

\section{Limitations}

The robustness evaluation uses a modest number of documents per attacker;
flip rates are directional rather than precise. The orthogonal-feature
baselines are corpus statistics estimated on this dataset and may not transfer
to other domains. P2/oracle numbers use test labels for calibration and are
explicitly upper bounds, not deployable scores; only the P1 number ($0.857$)
is leakage-free. The de-AI component is strictly a robustness probe; we do not
endorse using it to evade academic-integrity systems. Finally, our study is
English-only and single-domain at test time.

\section{Conclusion}

On a realistic distribution-shift detection benchmark, a frozen detector's
apparent failure (macro-F1 $0.62$) is almost entirely an operating-point
artifact: calibration plus an in-distribution threshold recovers $0.886$.
A region-aware two-stage rule that routes ambiguous-margin documents to an
orthogonal-feature classifier reaches a deployable $0.857$, beats the prior
best single model, and recovers the hardest human-like generator. The same
orthogonal stage is its main paraphrase-attack surface, and a hierarchical
capacity extension overfits and fails---together arguing that calibration and
distribution shift, not model capacity, are the levers that matter.

\bibliography{refs}
\bibliographystyle{acl_natbib}

\end{document}



